\section{Preliminaries}\label{section2}

For an algebra $R$ over a field $K$ a linear operator $\delta\colon R\rightarrow R$ is called a derivation if it satisfies the Leibniz law $\delta(ab) = \delta(a)b + a\delta(b)$.
The kernel of a derivation $\delta$ is denoted by~$R^{\delta}$ and the elements of the kernel are called $\delta$-constants (or just constants when this is not confusing).
A~derivation~$\delta$ is called locally nilpotent if for any $r \in R$ there exists a natural number~$n$ (which depends on~$r$) for which $\delta^n(r) = 0$. The function
\[
\deg(r) = \max\big(d\,|\, \delta^d(r) \neq 0\big), \qquad \deg(0) = -\infty,
\]
is a degree function with familiar properties:
\begin{gather*}
\deg(r_1r_2) = \deg(r_1) + \deg(r_2),\qquad \deg(r_1 + r_2) \leq \max(\deg(r_1), \deg(r_2)),\\
\deg(r_1 + r_2) =\max(\deg(r_1), \deg(r_2)) \qquad \text{when} \quad \deg(r_1) \neq \deg(r_2),\\
\deg(\delta(r)) = \deg(r) - 1\qquad \text{if} \quad \delta(r) \neq 0.
\end{gather*}

The set of all lnds (locally nilpotent derivations) of $R$ is denoted by $\text{LND}(R)$.

The intersection $\bigcap R^{\delta}$, $\delta \in \text{LND}(R)$, of kernels of all locally nilpotent derivations of $R$ is denoted by $\text{AK}(R)$ (absolute Konstanten of~$R$, sometimes denoted as $\text{ML}(R)$).

If $\delta \in \text{LND}(R)$ and characteristic of $K$ is zero then the linear operator
\[
\exp(\delta)=1+\frac{\delta}{1!}+\frac{\delta^2}{2!}+\cdots
\]
is an automorphism of $R$.

In the sequel we fix a field $K$ of characteristic 0 and consider the polynomial algebra $K[x,y]$ and the free associative algebra $K\langle X,Y\rangle$. Let
\[
\pi\colon \ K\langle X,Y\rangle\to K[x,y]
\]
be the natural homomorphism. We denote the elements $U$, $V$, etc.\ of $K\langle X,Y\rangle$ by upper case symbols and their images under $\pi$ by the same lower case symbols $u$, $v$, etc. Let $C$ be the commutator ideal of $K\langle X,Y\rangle$. It is generated by the commutator
\[
T_1 = [Y, X] = YX - XY.
\]

By the theorem of Jung--van der Kulk \cite{J, K}, the automorphisms of $K[x,y]$ are tame, i.e., are compositions of affine automorphisms
\[
x \rightarrow a_1 x + a_2 y + a_3, \qquad y \rightarrow b_1 x + b_2 y + b_3, \qquad a_i,b_i\in K,\qquad a_1b_2 - a_2b_1 \neq 0,
\]
and triangular automorphisms
\[
x \rightarrow x, \qquad y \rightarrow y + p(x),\qquad p(x)\in K[x].
\]

A similar theorem of Czerniakiewicz \cite{Cz1, Cz2} and Makar-Limanov \cite{ML} states that the automorphisms of $K\langle X,Y\rangle$ are also tame. Therefore
\[
\Psi(T_1) = cT_1, \qquad c \in K^{\ast},
\]
for any automorphism $\Psi$ of $K\langle X,Y\rangle$ (indeed, just check that this is true for affine and triangular automorphisms).

The structure of the automorphism groups of $K[x, y]$ and $K\langle X,Y\rangle$ is also known, it is a~free product of the subgroups of affine and triangular automorphisms with amalgamation along their intersection~\cite{Se}. So we can think that there is a group $H$ isomorphic to $\operatorname{Aut}K[x,y]$ and $\operatorname{Aut}K\langle X,Y\rangle$ which acts on $K[x, y]$ and $K\langle X,Y\rangle$.

Any automorphism of $K\langle X,Y\rangle$ induces an automorphism of $K[x, y]$ and, since the structure of the group $H$ insures that this is one to one correspondence, any automorphism of $K[x,y]$ can be uniquely lifted to an automorphism of $K\langle X,Y\rangle$.

We shall use below a lexicographic ordering of monomials of $K\langle X,Y\rangle$ defined by $Y\gg X > 1$ and denote by $\overline{S}$ the leading monomial of $S \in K\langle X,Y\rangle$.

In the sequel we shall show that we can reduce our considerations to the case when the lnd~$\Delta$ is such that
\[
\Delta(X)=0,\qquad \Delta(Y)=F=f(X),
\]
where $0\not=f(x)\in K[x]$. In this special case we shall define the operator $\boxdot$ on $K\langle X,Y\rangle$ by
\[
\boxdot(A)= YAF - FAY,\qquad A\in K\langle X,Y\rangle,
\]
and shall fix the sequence $T_1,T_2,\ldots$, starting with $T_1= YX - XY$ and then inductively
\[
T_{i+1} = \boxdot^i(T_1).
\]

